\subsection{ARCH-136-PD-ZEROFORM: a prime-uniform secondary-range rigidity}
\label{sec:arch136pdzeroform}

\paragraph{Scope and inherited input.}
This addresses a fixed-form-degree secondary PD operation suggested
by the cote 136 structure on physical p.~111, footer 110, with the
weight bounds of physical pp.~117 and 119, footers 116 and 118.
It excludes one proposed invariant for the common-source question
on physical p.~121, footer 120. The source files are respectively
\texttt{ARCHIVE\_136\_scan111.tex},
\texttt{ARCHIVE\_136\_scan117.tex},
\texttt{ARCHIVE\_136\_scan119.tex}, and
\texttt{ARCHIVE\_136\_scan121.tex}.

The inherited AQ-136-PD-RIGIDITY identifies primary divided-power
operations of each fixed order on positive even cohomology, with
a cutoff depending on that order. We do not reprove or rename it.
Here the inputs are cochains in degrees zero and one, the output
is in degree zero, and one cutoff works for \emph{all primes at
once}. The only inherited input is AQ-136-COMPARISON in degrees
zero and one; the additional arguments below have tag none.
A, M, and outside preprints are unused.

\paragraph{Statement.}
Let \(C\in\mathcal A_{\mathbf Z}\). There is \(W(C)\geq1\), independent
of the prime \(p\) and of the simplicial set \(X\), with the following
properties. Let \(w\geq W(C)\), and choose homogeneous cochains
\[
 a\in C^{w,1}(X),\qquad b\in C^{w,0}(X),\qquad
                         da=0,\quad db=p a .
 \tag{136.ZF1}
\]
Then \(a\gamma_{p-1}(b)\) is exact. If
\(du=a\gamma_{p-1}(b)\), define
\[
 \mathcal R_p(a,b)=
       [p u-\gamma_p(b)]
   \quad\hbox{in }H^{pw,0}(C(X))/pH^{pw,0}(C(X)).
 \tag{136.ZF2}
\]
This class is independent of \(u\), and natural under strict
differential PD maps when the displayed data are carried along.
Under the coefficient comparison, its full range over all pairs
(136.ZF1) is
\[
 \begin{cases}
  H^0(X;\mathbf F_p)\,T^{[pw]},& w=p^r\text{ for some }r\geq0,\\
  \{0\},&\text{otherwise}.
 \end{cases}
 \tag{136.ZF3}
\]
In particular this range is independent of the admissible theory,
uniformly in \(p\), once the input weight is at least \(W(C)\).

\paragraph{Proof.}
Choose \(W(C)\geq1\) so that AQ-136-COMPARISON identifies
\[
 H^{v,j}(C(X))\simeq H^j(X;\mathbf Z)\,T^{[v]}
               \quad(j=0,1;\ v\geq W(C))
 \tag{136.ZF4}
\]
for every \(X\). No prime enters this choice.
Integral \(H^1(X;\mathbf Z)\) is torsion-free, for arbitrary \(X\).
For an elementary verification of the needed assertion, if
\(p\alpha=\delta f\) for an integral one-cocycle, the values of
the zero-cochain \(f\) are congruent modulo \(p\) on every connected
component. Subtract a constant integer lift of that residue on
each component and divide by \(p\); the resulting zero-cochain
has coboundary \(\alpha\).

Equation (136.ZF1) and (136.ZF4) therefore imply \(a=dc\) for some
homogeneous \(c\in C^{w,0}(X)\). Put
\[
                              s=b-pc .
\]
Then \(ds=0\). As there are no negative cochain degrees, the
degree-zero comparison says that \(s\) is an actual coefficient
section: a locally constant integer function on \(X\) times
\(T^{[w]}\). All of \(b,c,s\) have positive weight, so their
divided powers are defined.

Use the PD addition and scaling identities to define
\[
 u_0=\sum_{i=1}^p
       p^{i-1}\gamma_i(c)\gamma_{p-i}(s).
 \tag{136.ZF5}
\]
Since \(ds=0\) and \(dc=a\), direct differentiation gives
\[
 \begin{split}
 du_0
 &=a\sum_{i=1}^p
       p^{i-1}\gamma_{i-1}(c)\gamma_{p-i}(s)\\
 &=a\gamma_{p-1}(pc+s)
  =a\gamma_{p-1}(b),\\
 p u_0-\gamma_p(b)&=-\gamma_p(s).
 \end{split}
 \tag{136.ZF6}
\]
This proves existence of \(u\) and the asserted residue formula.
Changing \(u\) adds \(p\) times a closed zero-cochain, and hence
does not change (136.ZF2). The differential of its representative
is \(p a\gamma_{p-1}(b)-\gamma_{p-1}(b)\,p a=0\).
Naturality follows from the actual formulas. If the chosen
primitive \(c\) of \(a\) is changed, \(s\) changes by \(p\) times
a closed zero-cochain. The PD binomial expansion shows that
\(\gamma_p(s)\) modulo \(pH^0\) is unchanged, in agreement with
the already proved independence from \(u\).

On a connected component write \(s=mT^{[w]}\), with \(m\in\mathbf Z\).
The coefficient identity is
\[
 \gamma_p(s)=m^p A_{p,w}T^{[pw]},
 \qquad
 A_{p,w}=\frac{(pw)!}{p!(w!)^p}\in\mathbf Z .
 \tag{136.ZF7}
\]
Integrality follows either from iterated divided powers or by
counting partitions into \(p\) unordered blocks of size \(w\).
Let \(s_p(w)\) be the sum of the base-\(p\) digits of \(w\).
Counting multiples of \(p,p^2,\ldots\) in each factorial gives
\[
 \begin{split}
 v_p(A_{p,w})
 &=\frac{pw-s_p(pw)}{p-1}-1
       -p\,\frac{w-s_p(w)}{p-1}\\
 &=s_p(w)-1 .
 \end{split}
 \tag{136.ZF8}
\]
This vanishes exactly when \(w\) is a power of \(p\).
In that case \(A_{p,w}\) is a unit modulo \(p\), and
\(m\mapsto m^p\) is the identity modulo \(p\). Thus the values
of \(-\gamma_p(s)\) fill the whole rank-one coefficient group
modulo \(p\); otherwise they are all zero. The argument applies
independently on all connected components, including an infinite
set of components. Every such value occurs already for the
pair \(a=0,b=s\), with \(u=0\). This proves both inclusions in
the range description (136.ZF3).

Finally the output weight \(pw\) is at least \(w\), so it remains
in the same range (136.ZF4). There was no order-dependent
cohomological bound anywhere in the argument. This proves the
prime-uniform assertion. \(\square\)

\paragraph{Exact consequence and limitation.}
For two admissible theories take the maximum of their two bounds.
The ranges (136.ZF3) then agree for every prime and every
simplicial set. Thus this particular zero-form secondary range
cannot separate components, even by taking an unbounded family
of primes with one input-weight bound. We have not asserted
that every system of cochain primitives can be transported along
an arbitrary arrow, or that all secondary PD operations are of
this form. The integral common-source and terminality questions
remain unaffected by that stronger, unjustified inference.
